\documentclass{article}

\begin{document}
\title{HOPE draft}


\section{Motivation}
We show how to combine the approaches to algebraic effects, together with failures of pattern matching.
We do our development in 3 phases, we show that 

Effect monad ala Kiselyov satisfies the monad laws.
Effect monad ala Kiselyov satisfies the graded monad laws.

\section{Background}

\subsection{Algebraic effects as monads and graded monads}

\subsection{Handling effects}

\section{Our results}

We have produced an intermediate construction between the variations of Kiselyov in Agda.
TODO add short explanation of similarities and differences

\section{Future directions}
As much we have constructed the monad in agda and verified the various laws. For practical purposes there needs to be more than just a monadic structure.
We want to reason about the local code using the eponymous algebraic laws. We would like to attach equations the handlers should satisfy to the monad.
There are multiple interesting directions to take this in. 
We would also like to use Quotient inductive types to provide a view of operations for reasoning and correct-by-construction definition of handlers,
which abstracts away from the concrete choices when reasoning, yet forces us to prove satisfiability of the algebriac equations by some handler.
To extend the proofs from our current work to the quotient case we need to prove a type-theoretic formulation of algebraic freedom \cite{ncatlab}.
To verify that all the constructions preserve all equations. We note that by theorem 3.1 and 3.2 \cite{ncatlab} this result should follow in Homotopy type theory.
The reason for this being that we need the ambient category to be locally small and complete with homotopy type theory we should be able to get all colimits.
\end{document}